\subsection{From the LP master to certified integer endpoints}
\label{sec:branch-price}

Column generation certifies the Dantzig--Wolfe relaxation, but aggregate
frontiers are defined over event \emph{partitions}, not fractional mixtures of
partitions.  This distinction is substantive: Appendix~\ref{app:integer-bp}
gives an instance with
$\operatorname{OPT}_{\mathrm{LP}}=4$ and
$\operatorname{OPT}_{\mathrm{IP}}=3$.  Solving a restricted integer master after
root-node column generation is therefore only a lower bound.

We close the integer problem by branch-and-price.  At each node, the restricted
master is solved to LP optimality and the rooted interval oracle searches for
negative-reduced-cost run columns under the node decisions.  We first branch on
an optional row whose aggregate usage is fractional.  Once every row usage is
integral, we use Ryan--Foster branching on a pair of rows with fractional
co-membership: one child requires the pair to occur together in every selected
run, and the other forbids joint membership.  Together/separate decisions are
expanded into forced-in/forced-out pricing cases.  Each case adds only unit
fixings to the same fixed-span interval LP, so node pricing remains exact and
integral.

\begin{theorem}[Finite exact branch-and-price]
\label{thm:bp}
For a finite candidate-row set, suppose every node master is solved exactly and
every admissible improving run column is found by the rooted interval pricing
oracle.  Buffer-usage branching followed by Ryan--Foster pair branching
terminates after finitely many nodes and returns the optimal integer ordered-run
partition, or certifies infeasibility.
\end{theorem}

The proof uses two facts.  Every branch is a disjunction covering the parent
worlds, and the family of possible run-member masks is finite.  Moreover, if
all row usages and all pair co-memberships are integral, every positive-weight
master column induces the same equivalence relation on the used rows; after
coalescing identical masks, the node solution is an integer partition.  The
full argument and certificate checks are in Appendix~\ref{app:integer-bp}.
The theorem is an exactness result, not a polynomial-time claim: the number of
pricing cases and branch nodes can grow exponentially.

\begin{table}[t]
\centering
\caption{Exact-time NYC integer audit.  ``All columns'' is used only as an
independent small-instance verifier; branch-and-price itself generates columns
on demand.}
\label{tab:bp}
\small
\begin{tabular}{@{}rrrrr@{}}
\toprule
$C$ & All columns & IP optimum & B\&P nodes & Time (s) \\
\midrule
2 & 54    & 4  & 1 & 0.74 \\
3 & 394   & 8  & 4 & 10.47 \\
4 & 1,719 & 12 & 2 & 9.31 \\
\bottomrule
\end{tabular}
\end{table}

On the audited cohort, branch-and-price agrees with the independently enumerated
integer master for all three capacities (Table~\ref{tab:bp}).  On the locked
nonintegral witness, it processes 17 nodes and closes the gap from four to
three.  These experiments establish an exact medium-instance engine and expose
its scaling variables---run-column density, root relaxation gap, pricing cases,
and branch-node count---rather than claiming city-wide relation recovery.
